\subsection{7. Discussion}\label{discussion}

\hypertarget{interpretation-of-results}{%
	\subsubsection{7.1 Interpretation of Results}\label{interpretation-of-results}}

A2C outperforms Double DQN on moving IoT service composition while building upon the STR-based selection from prior work \cite{neiat2021}. Three factors drive improvements.

First, actor-critic gives stable learning by combining value estimation
with direct policy optimization. The advantage function reduces variance
while keeping updates unbiased, so the agent learns effectively from
fewer samples.

Second, the relative polar coordinate representation captures spatial
relationships efficiently, enabling the agent to learn effective service
selection without explicit velocity or trajectory prediction.

Third, separate networks specialize processing for spatio-temporal
states. More parameters and training time, but better performance when
movement patterns matter.

The STR-based selection continues driving decisions---the agent learns
to maximize expected capacity while maintaining stability.

\hypertarget{comparison-with-prior-work}{%
	\subsubsection{7.2 Comparison with Prior
		Work}\label{comparison-with-prior-work}}

A2C improves over Double DQN on both datasets.

ATC (Indoor): A2C Separate achieves 95.2\% versus 92.4\% at low
mobility (2.8\% improvement). At high mobility (10 km/h), the gap grows
to 13.9\% (85.7\% vs 71.8\%).

Illinois (Outdoor): At high mobility (10 km/h), A2C Separate reaches
80.3\% versus 64.2\% for Double DQN---a 16.1\% absolute improvement.
Outdoor trajectories with more varied movement patterns benefit from
the learned policy.

Capacity satisfaction improves because A2C better captures spatial
relationships and selects services with better capacity margins.

\hypertarget{practical-implications}{%
	#### 7.3 Practical Implications}\label{practical-implications}}

\textbf{Edge Deployment}: A2C with shared networks balances performance
and computation for edge deployment. Inference cost allows real-time
decisions on edge devices.

\textbf{Stability}: A2C re-composes less frequently than Double DQN,
reducing service disruption and overhead for production systems.

\textbf{STR Quality}: STR-derived capacity provides a grounded service
quality metric tied to spatial proximity.

\hypertarget{limitations}
#### 7.4 Limitations}\label{limitations}

Three limitations point to future work:

\textbf{STR Model}: Uses simplified distance-based models. More complex
propagation---multipath fading, shadowing, interference---needs
investigation.

\textbf{Centralized}: Assumes centralized composition. Distributed
multi-agent approaches could scale better for large IoT systems.

\textbf{Validation}: We tested on real GPS from ATC and Illinois
datasets, but service availability is simulated. Future work should
use real IoT service availability data.

\begin{center}\rule{0.5\linewidth}{0.5pt}\end{center}

\hypertarget{conclusion}{%
